% ============================================================================
% Section 11: Conclusion
% ============================================================================

\section{Conclusion}
\label{sec:conclusion}

Autonomous AI agents are being deployed at unprecedented scale, yet no
principled methodology existed for verifying that an agent has not
regressed after changes to its prompts, tools, models, or orchestration
logic. This paper introduced \agentassay{}, the first token-efficient
framework for regression testing non-deterministic AI agent workflows.

Our key insight is that agent testing requires a paradigm shift from
binary verdicts to \emph{stochastic, three-valued} verdicts
(\PASS{}/\FAIL{}/\INCONC{}) grounded in statistical hypothesis testing.
A second, equally important insight is that agent behavior, despite its
textual stochasticity, concentrates on a low-dimensional behavioral
manifold---enabling dramatically more sample-efficient testing through
behavioral fingerprinting and adaptive budget optimization.
Building on these foundations, we made ten contributions:

\begin{enumerate}[leftmargin=*]
  \item \textbf{Stochastic test semantics} with the $(\alpha, \beta, n)$-test
    triple, three-valued verdict function, and provably sound regression
    detection (\cref{thm:soundness,thm:power}).

  \item \textbf{Agent coverage metrics}---a five-dimensional tuple
    $(C_{\text{tool}}, C_{\text{path}}, C_{\text{state}},
    C_{\text{boundary}}, C_{\text{model}})$---that measure test
    thoroughness for stochastic agent systems (\cref{thm:coverage-mono}).

  \item \textbf{Agent mutation testing} with four classes of
    domain-specific operators and a formal mutation adequacy theorem
    (\cref{thm:mutation-adequacy}).

  \item \textbf{Metamorphic relations} tailored to agent workflows,
    addressing the oracle problem for open-ended agent tasks.

  \item \textbf{CI/CD deployment gates} as statistical decision
    procedures with configurable risk thresholds.

  \item \textbf{Contract integration} with the \agentassert{} framework,
    using behavioral contracts as formal test oracles with provable
    correspondence (\cref{prop:verdict-contract}).

  \item \textbf{SPRT adaptation} for cost-efficient agent testing,
    achieving \textbf{78\%} trial savings while maintaining error
    guarantees (\cref{thm:sprt}).

  \item \textbf{Behavioral fingerprinting} that maps execution
    traces to compact vectors on a low-dimensional manifold,
    enabling multivariate regression detection with provably
    higher power per sample (\cref{thm:fingerprint-regression}).

  \item \textbf{Adaptive budget optimization} that calibrates trial
    counts to actual behavioral variance, achieving 4--7$\times$
    reduction for stable agents (\cref{thm:adaptive-budget}).

  \item \textbf{Trace-first offline analysis} that eliminates live
    agent executions for four of six test types with formal soundness
    guarantees, enabling zero-cost coverage, contract, and metamorphic
    testing on production traces (\cref{thm:trace-first}).
\end{enumerate}

The combined token-efficient testing framework achieves
5--20$\times$ cost reduction (\cref{thm:combined-efficiency}), bringing
rigorous statistical agent testing within CI/CD budget constraints for
the first time.

Experiments across \textbf{5} models, \textbf{3} scenarios, and
\textbf{7{,}605} trials (\$227 total cost) demonstrated that \agentassay{}
provides statistically rigorous regression detection with practical
efficiency. Behavioral fingerprinting achieved \textbf{86\%} detection
power where binary pass/fail testing had \textbf{0\%}, SPRT reduced
trial counts by \textbf{78\%} consistently across all three scenarios,
and the full token-efficient pipeline achieved \textbf{100\%} cost
savings through trace-first offline analysis---all while maintaining
identical $(\alpha, \beta)$ guarantees.

\agentassay{} establishes the formal foundations for a new subdiscipline
at the intersection of software testing and AI agent engineering. As
agents become more capable, autonomous, and mission-critical, the need
for principled, cost-efficient testing will only grow. We believe that
the concepts introduced in this paper---stochastic verdicts, agent
coverage, agent mutation testing, behavioral fingerprinting, and
token-efficient testing---will become standard practice in the emerging
field of \emph{Agent Software Engineering}.

\paragraph{Availability.}
\agentassay{} is available as open-source software under the
Apache~2.0 license. The implementation comprises ${\sim}$20{,}000 lines of
Python with 751 tests and adapters for 10 agent frameworks. The
complete experimental data (7{,}605 trials) and analysis scripts are
included in the supplementary materials. Zenodo DOI:
\href{https://doi.org/10.5281/zenodo.18842011}{10.5281/zenodo.18842011}.
